\chapter{Algebraic tools for singular contacts}
\label{ch:mechanisms}
A singular contact has vanishing quadratic terms. The following identities
express its higher-order terms using roots with a positive residue Gram.
Chapter~\ref{ch:tensor} then turns that Gram into an exact sum of squares.
\section{Finite algebra and contact completion}
\paragraph{Finite penalty lemma.}
Let $R$ be a real symmetric quadratic form and $L$ a real linear
map.  If $R$ is positive definite on $\ker L$, then
$R+M\|Lz\|^2$ is positive definite for some real $M>0$.

\emph{Proof.}
Choose a linear complement of $\ker L$.  The restriction of $L$
to that complement has a positive least singular value.  In these
coordinates the kernel block of $R$ is positive definite, and
adding $M L^TL$ makes the other Schur complement positive for
sufficiently large $M$. \hfill$\square$

For an opening argument one must additionally exhibit a bounded
remainder after subtracting the divergent squares, and a bounded
right inverse for their independent hard rows.  If a family has a
uniform positive lower bound before subtraction, its limit on the
limiting kernel is positive: lift a kernel vector to exact kernels
with that right inverse and take the real limit.  The finite penalty
lemma then supplies finite caps.  The applications in Appendices B and C give the
required bounded minimizers and independent rows.

\paragraph{Small-cost helper lemma.}
Suppose the main roots have a strict real residue Gram and an exact
identity has the form $A(z,H)-H^2=\sum R_i^2$, where every
coefficient of the quadratic form $A$ on the main roots $z$ and
$H$ is infinitesimal.  Then any prescribed finite positive diagonal
budget $cH^2$ can be introduced by the exact replacement
\[
          F=(F-cA+cH^2)+c\sum R_i^2.
\]
It changes the main residue Gram by zero and retains the displayed
positive squares.  Finitely
many helper roots can be treated in a stated order, after which a
real Schur complement handles their bounded cross coefficients.
When the identities are triangular, introduce the auxiliary roots in
that order. In the double-infinitesimal C1 case this means obtaining
the radial pivot before the angular roots.


\section{Finite polynomial division at a contact}
\label{component:pfabc29e28ab3}
% Optional module: the source notation is that of (UZ4)--(UZ6).
% All statements below are exact; no localized ring is used.


Write $d=\delta$ in this section; this local parameter is distinct from
the global Bell coefficient $d_{\rm B}$. Put
\[
 R_0=\mathbb Q[d,k,b,a,t],\qquad R=R_0[u,w,v],\qquad
 J=(u-k,w-dk,v-dk^2b^2).
\]
For $F\in R$, write $F_c=F(k,dk,dk^2b^2)$. Define
\[
 A=wv,\quad B=dv,\quad C=db^2uw,\quad D_0=d^2k^3b^2,
 \qquad M=(A,B,C),
\]
\[
 x=B-C,\qquad y=A-kC,\qquad z=uB-kC,\qquad h=kC-D_0.
\]
Equivalently, write $g_0=x$, $h_{13}=y$, $h_{23}=z$, and $h_{34}=h$. Then
$D_0=a_4^2$, whereas the source monomial in $g_4$ is $kD_0$.
Adjoin the separated radicals used in the four-monomial identity,
\[
 a_1=\sqrt{wv},\quad a_2=\sqrt{duv},\quad
 a_3=b\sqrt{dkuw},\quad a_4=dk^{3/2}b,
 \qquad q=a_1a_2-a_3a_4.
\]
Only $q^2\in R$ is needed in the scalar lift. The following three identities
are the entire conductor calculation:
\begin{align}
 C(u-k)&=z-ux,\label{pfabc29e28ab3:CC-u}\\
 C(w-dk)&=dy-wx,\label{pfabc29e28ab3:CC-w}\\
 C(v-dk^2b^2)&=b^2(wz+dky-kwx).\label{pfabc29e28ab3:CC-v}
\end{align}
Thus $CJ\subset(x,y,z)$ without inverting any parameter.
For
\[
 r_*=y+z-h-2kx
\]
we also have the two exact identities
\begin{align}
 r_*&=(v-dk^2b^2)\bigl((w-dk)+d(u-k)\bigr)
       -dkb^2(u-k)(w-dk),\label{pfabc29e28ab3:CC-star}\\
 kCr_*&=q^2-yz.\label{pfabc29e28ab3:CC-repair}
\end{align}
The first says $r_*\in J^2$. The second supplies the otherwise missing
conductor factor through one of the original circuit squares.

\paragraph{Explicit monic division.}
Every $S\in R$ has a polynomial decomposition
\[
 S=S_c+(u-k)T_u+(w-dk)T_w+(v-dk^2b^2)T_v.                 \tag{CC1}
\]
One obtains it by substituting $u=k$, then $w=dk$, then $v=dk^2b^2$,
and using $X^n-Y^n=(X-Y)\sum_{j=0}^{n-1}X^{n-1-j}Y^j$ at each step.
There are no parameter denominators. If $S_c=0$, equations
\eqref{pfabc29e28ab3:CC-u}--\eqref{pfabc29e28ab3:CC-v} give
\[
 CS=x(-uT_u-wT_w-b^2kwT_v)
     +y(dT_w+dkb^2T_v)+z(T_u+b^2wT_v).                 \tag{CC2}
\]

\begin{lemma}[First-order and restricted second-order contact lifts]
In the above ring, the following assertions hold, with polynomial
coefficients in $R$:
\[
 M\cap J\subset (x,y,z),\qquad
 \bigl(M^2+(D_0^2)\bigr)\cap J^2
       \subset (x,y,z,h,q)^2.                         \tag{CC3}
\]
Here $(D_0^2)$ is the principal ideal generated by $D_0^2$; the second
assertion does not replace it by all mixed products involving $D_0$.
\end{lemma}
\begin{proof}
For the first assertion, write $F=R_AA+R_BB+R_CC$. Since
$A=y+kC$ and $B=x+C$,
\[
 F=R_Ay+R_Bx+CS,\qquad S=kR_A+R_B+R_C.
\]
The contact value is $F_c=C_cS_c$, where $C_c=d^2k^2b^2$ is a nonzero
polynomial. Since $R_0$ is a domain, $F_c=0$ implies $S_c=0$.
Apply (CC2). This cancellation proves a zero identity; it introduces no
coefficient $1/C_c$.

For the second assertion choose any displayed representation
$F=\sum_{i\le j}f_{ij}m_im_j+\rho D_0^2$, with
$(m_1,m_2,m_3)=(A,B,C)$. Substitute
\[
 A=y+kC,\qquad B=x+C,\qquad D_0=kC-h.
\]
This gives
\[
 F=Q_2+C(xH_x+yH_y+zH_z+hH_h)+C^2S,\qquad H_h=-2k\rho,       \tag{CC4}
\]
where $Q_2\in(x,y,h)^2$ and initially $H_z=0$.
The contact value again implies $S_c=0$. Use (CC2) on $CS$ and absorb
$C^2S$ into the four coefficients in (CC4). This leaves $H_h=-2k\rho$
unchanged and gives the same formula without $C^2S$.

The first derivatives of $F\in J^2$ vanish at contact. The first three
root gradients have determinant
\[
 \det\left(\frac{\partial(x,y,z)}{\partial(u,w,v)}\right)_c
       =-d^4k^4b^4\ne0.                              \tag{CC5}
\]
Equation~\eqref{pfabc29e28ab3:CC-star} gives a gradient relation with coefficient vector
$(-2k,1,1,-1)$. Its coefficient at $h$ is $-1$.
Set $\lambda=2k\rho_c$. Subtract
$\lambda(-2k,1,1,-1)$ from $(H_x,H_y,H_z,H_h)$.
The last resulting contact value is zero, and the first three satisfy the
homogeneous linear system with determinant (CC5). Multiplication by its
adjugate and cancellation in the domain $R_0$ show that those three
contact values are zero too. No inverse of the determinant is used in a
coefficient: each resulting polynomial $J_i$ lies in $J$ and is divided
only by the monic procedure (CC1).
Consequently
\[
 F=Q_2+2k\rho_cCr_*+C\sum_{i=1}^4r_iJ_i
   =Q_2+2\rho_c(q^2-yz)+\sum_{i=1}^4r_i(CJ_i),          \tag{CC6}
\]
where $(r_1,r_2,r_3,r_4)=(x,y,z,h)$. Each $CJ_i$ is explicitly a
combination of $x,y,z$ by (CC2), which proves the assertion in the
unlocalized polynomial ring.
\end{proof}

The explicit scalar and both state-normal coefficients in
Section~\ref{sec:c1-bounded-reciprocal} satisfy these ideal conditions.
There the same monic divisions also preserve a positive coefficient
weight. This supplies bounded corrections with zero residue, which
can be added to the strict Gram matrix.

\section{A real two-block allocation}
\label{component:p6b22e944bd7c}
\paragraph{A rank-one allocation lemma.}
Let $P$ be a positive definite real symmetric matrix, let $\ell,b$ be
real columns, and let $H,c_0>0$. Set
\[
 d=\ell^TP^{-1}\ell,\qquad e=\ell^TP^{-1}b,\qquad
 h=c_2-b^TP^{-1}b.
\]
Suppose
\[
 p(T)=c_0+(h+4Hdc_0)T+4H(dh+e^2)T^2>0\quad(T\ge0),
 \qquad dh+e^2>0.                                      \tag{UA1}
\]
There exists a real $\lambda$ such that
\[
 |\lambda|<\sqrt{4Hc_0},\qquad
 h+2e\lambda-d\lambda^2>0.                             \tag{UA2}
\]
Indeed put $s=\sqrt{4Hc_0}$. If $d>0$ and $|e|<ds$, take
$\lambda=e/d$. If $|e|\ge ds$, the value at the appropriate endpoint
is $h+2|e|s-ds^2$. Were this nonpositive, writing
$x=|e|-ds\ge0$ would give
$h+ds^2\le-2sx$ and $dh+e^2\le x^2$; the quadratic in (UA1)
would then have a nonpositive value at a positive argument. Thus a
nearby interior $\lambda$ works. This includes equality $|e|=ds$.
The scalar interval argument also holds in a limiting case $d=0$;
then $e\ne0$ follows from $dh+e^2>0$. (For a finite positive $P$
itself, $d=0$ would force $\ell=e=0$, so this case occurs only
when using the continuous scalar extension at a boundary.)

Consequently the two real matrices
\[
 \begin{pmatrix}P&b-\lambda\ell\\(b-\lambda\ell)^T&c_2\end{pmatrix},
 \qquad
 \begin{pmatrix}4H&\lambda\\\lambda&c_0\end{pmatrix}       \tag{UA3}
\]
are positive definite. This is an elementary real lemma. Its inverse
and square root are never taken in the tensor ring.


\section{Noncommuting testing residues}
\label{component:pe25c50379365}
\label{component:p76173e99815f}
\label{component:p24f8a75e651e}
\label{component:p4e9ae18b7af4}
\label{component:pa2cbe6ecbeb4}
\label{component:pb16e5066334b}
\label{component:pf47f6a4e3a60}
\label{component:p94dfbc8a5970}
\label{component:p2c3fc22497a5}
% Standalone replacement for the NONCOMMUTING portions of components4--12.
% Uses only the tensor foundation (positive residue, strict bounded Gram,
% pure-factor cancellation), standard real algebra and the displayed parent.

\label{sec:noncommuting-refined}

Use the two labelled fields and the bounded tensor ring of
Chapter~\ref{ch:tensor}. The strict-Gram lemma applies to arbitrary
root columns, including unbounded rows.

\subsection{Parent, physicality and contact curvature}

For real homogeneous state coordinates $\psi=(W,X,Y,Z)$ and half angles
$u,v$, set
\[
 D_u=1+u^2,\quad D_v=1+v^2,\quad
 M=vW-X+uvY-uZ,\quad J=uW+uvX-Y-vZ.
\]
All nine source coefficients $d,\lambda,n,m,c,A,D,\alpha,\beta$ are
nonnegative. The entire parent is
\begin{align}
 F={}&4\lambda D_u(W+vX)^2+4nD_uD_vX^2+4mM^2
       +4cD_v(uW-Y)^2+4AD_uD_vY^2\notag\\
 &+(4D+2d)J^2+(4\alpha+2d)D_u(vY-Z)^2
       +(4\beta+2d)D_v(X+uZ)^2-2dD_uD_vZ^2.
 \tag{NC1}
\end{align}
Write $C$ for the coefficient of $d$ in this expression. A source is
\emph{physical} if its state matrix is positive semidefinite at every
measurement pair over its own field $K_c$. The claim below concerns
evaluation of (NC1) at an independent test over $K_t$.

\begin{theorem}\label{thm:noncommuting-refined}
For every physical source in (NC1), its tensor evaluation is a finite
sum of squares at every independent test for which both measurement
pairs are noncommuting in residue, and every homogeneous real or
complex state. No relation between source scales or between source
and testing scales is assumed.
\end{theorem}

Local frame changes and a projective state chart are pure-testing
operations. At a noncommuting residue one may take $u,v$ to be bounded
units, and normalize one state coordinate to a bounded unit. The
original slack is restored by positive pure-testing norm factors.
For complex states, apply the real result separately to real and
imaginary parts, since the state matrix is real symmetric. We prove
the real case.

If $d=0$, the parent is already a sum of event squares.
For the remaining $d>0$ case, use the eight outer coefficients
$q_{ij}$ and the positive quantities $t_i$ in (SG5)--(SG6), put
$q_{11}=h$, and denote the event coefficient by $c$ here. Equivalently
\[
 \det P=16D_u^2D_v^2\sum_{i,j=0}^2q_{ij}u^{2i}v^{2j}.
 \tag{NC2}
\]
Thus every coefficient in the constrained determinant calculations
below is specified by that displayed finite dictionary.

For $U=u^2,V=v^2>0$, let
\[
 L(U,V)=q_{11}+\sum_{(i,j)\ne(1,1)}q_{ij}U^{i-1}V^{j-1}.
 \tag{NC4}
\]
Every nonconstant coefficient is nonnegative. Over the real residue
field, a stationary zero of such a Laurent polynomial satisfies
\[
 \sum_{e\in E}w_e e=0,\qquad
 \nabla^2_{\log U,\log V}L=\sum_{e\in E}w_e ee^T,
 \quad w_e=q_eU^{e_1}V^{e_2}>0.                       \tag{NC5}
\]
Consequently its angular Hessian is positive definite exactly when
the active exponent vectors span $\mathbb R^2$. Logarithms are used
only for this real calculation; the displayed formulas are algebraic.

When $d>0$ and $uv\ne0$, the core restricted to $Z=0$ is
\[
 C|_{Z=0}=2\{(uW+uvX-Y)^2+D_uv^2Y^2+D_vX^2\},       \tag{NC6}
\]
which is positive definite in $(W,X,Y)$. Thus a physical state matrix
at a nonzero kernel has exactly three positive transverse directions.
The same observation supplies two or one positive transverse
directions after imposing one or two independent state-linear event
constraints, respectively.

\subsection{Exact local lifting with unequal hard weights}

We record the local argument that will turn real curvature into an
exact tensor certificate. Suppose, in a real coordinate neighborhood,
the complete physical parent has the form
\[
 F(x,y)=\sum_{i=1}^k a_i b_i(y)x_i^2+\sigma G(x,y),
 \qquad \varepsilon_i=\sigma/a_i>0\text{ infinitesimal},
 \tag{NC7}
\]
where $a_i,\sigma$ are pure-source, all coefficients of $b_i,G$ are
bounded, $\bar x=0$, $\bar y=y_0$, and $\bar b_i(y_0)>0$.
The case $k=0$ is allowed. Rational expressions are also allowed when
their denominators are products of pure-source units and real testing
polynomials nonzero at the residue. Multiply first by their squares;
the resulting polynomial has the same physical sign and the same
strictness at a zero, and the factors are canceled afterward.

Let $g(y)=\overline G(0,y)$. Physicality on the exact hard kernel
implies $g\ge0$ in its real neighborhood. If $g(y_0)>0$, the roots
$(x_i/\sqrt{\varepsilon_i},1)$ give an exact bounded Gram with residue
$\operatorname{diag}(\bar b_i(y_0),g(y_0))\succ0$:
every term of $G-G(0,y)$ containing $x_i$ acquires a factor
$\sqrt{\varepsilon_i}$ in this frame.

If $g(y_0)=0$ and $\nabla^2g(y_0)\succ0$, physicality also gives
$\nabla g(y_0)=0$. Put $x_i=\varepsilon_i s_i$ in the stationary
equations for $F/\sigma$ and divide the $x_i$ equations by their
nonzero factors. The exact equations are
\[
 2b_i(y)s_i+G_{x_i}(\varepsilon s,y)=0,\qquad
 \sum_i\varepsilon_i s_i^2\nabla b_i(y)
                       +G_y(\varepsilon s,y)=0.       \tag{NC8}
\]
At residue they have a block upper triangular Jacobian with diagonal
blocks $2\operatorname{diag}(\bar b_i(y_0))$ and
$\nabla^2g(y_0)$. The real implicit-function theorem gives a unique
local solution in fixed semialgebraic neighborhoods. Existence,
uniqueness and continuity bounds in those neighborhoods are fixed
first-order statements and transfer to $K_c$. They give an exact
source-field stationary point $(x_c,y_c)$ with
$x_{ci}=\varepsilon_i s_{ci}$ and bounded $s_c$, of the prescribed
residue. No center depends on a testing infinitesimal.

Finite polynomial Taylor expansion at this center, using the rows
\[
 h=\big((x_i-x_{ci})/\sqrt{\varepsilon_i}\big)_i
                    \mathbin{\oplus}(y-y_c),
\]
gives exactly
\[
 F/\sigma=a+h^TQh,\qquad
 a=F(x_c,y_c)/\sigma\ge0,\qquad
 \bar Q=\operatorname{diag}
   (\bar b_i(y_0),\tfrac12\nabla^2g(y_0))\succ0.      \tag{NC9}
\]
Indeed two hard deviations are allocated to their two scaled rows;
one hard deviation and a tangent deviation leave
$\sqrt{\varepsilon_i}$ in the coefficient; a hard packet's pure
tangent term leaves $\varepsilon_i$, because
$x_{ci}/\varepsilon_i$ is bounded. The $G$ terms have the same or
smaller factors. All higher Taylor terms leave an infinitesimal
deviation in the coefficient. The exact linear terms vanish at the
center. This proves the full-parent certificate, including all state
normals and higher terms. In particular, $k=0$ proves the ordinary
regular-contact case.

\subsection{Geometry of the event constraints}

The losing events form one bipartite cycle of local outcome rays,
in edge order
\[
 (\lambda,D,c,A,\alpha,m,\beta,n).                   \tag{NC10}
\]
For completeness its consecutive vertices may be taken as
$A_0^+,B_1^+,A_1^-,B_0^+,A_0^-,B_1^-,A_1^+,B_0^-$.
The four rays at either party are distinct at noncommuting residue
angles. CHSH-preserving relabelings act on $(x,y,a,b)$ as
\[
\begin{split}
g_A:(x,y,a,b)&\mapsto(1-x,y,a,(-1)^yb),\\
g_B:(x,y,a,b)&\mapsto(x,1-y,(-1)^{1-x}a,b),\\
g_T:(x,y,a,b)&\mapsto(y,x,(-1)^yb,(-1)^xa).
\end{split}
\]
They act on the edge indices of (NC10) by $i\mapsto5-i,3-i,2-i$
modulo eight. Hence they generate a cycle rotation and reflection.
They preserve the parent interpretation; chart changes contribute
only positive pure-testing homogeneous factors.

Product residue states need no further argument. For a product state
the CHSH expression satisfies
\[
 S=a_0(b_0-b_1)+a_1(b_0+b_1)
 \le |b_0-b_1|+|b_0+b_1|=2\max(|b_0|,|b_1|)\le2.
\]
The inequality is strict when both measurement pairs are
noncommuting. If a Bob modulus is one, the other is strictly less
than one, and equality would require Alice to be a common eigenstate
of her two noncommuting reflections. Thus $C$ has positive testing
residue and its actual pure-testing value is positive. The term $dC$
and all event packets are then sums of squares using separate roots.

For an entangled real residue state, write its amplitudes as an
invertible two-by-two matrix $\Psi$. A fixed Alice ray $a$ has
exactly one Bob projective ray satisfying $a^T\Psi b=0$, and
conversely. The vanishing events therefore share no endpoint in
(NC10): they form a matching and there are at most four of them.
Any three matching event vectors are state-linearly independent.
To see this, the span of two rank-one matrices with independent
left and right factors contains no other rank-one projective point;
three dependent rank-one matrices must therefore share a left or a
right factor, contrary to the matching condition. For $k\le3$, the
event amplitudes are consequently exact local projective-state
coordinates. Their coordinate inverse has only pure-testing unit
denominators and their packets have the diagonal form (NC7).

There are two perfect matchings. For $\{n,m,A,D\}$ their equations
are $X=Y=0$, $vW=uZ$, $uW=vZ$. Thus $u^2=v^2$; on either sign
branch $W\ne0$ and, after a sign gauge, $W=Z=1$, $v=u=t\ne0$.
With $t=u$ and $\delta=v-u$, the independent event-normal differentials are
\[
 dX,\quad dY,\quad
 dR_m=d\delta-dX+t^2dY-t\,dZ,\quad
 dR_D=t^2dX-dY-d\delta-t\,dZ.
\]
Their rank is four; the remaining tangent is $t$. The core is zero
on the curve and has differential
\[
 dC=2t\{(1+t^2)dX-(1+t^2)dY-dR_m+dR_D\}.           \tag{NC11}
\]
Every coefficient here is nonzero. The other perfect matching is
$\{\lambda,c,\alpha,\beta\}$; with $W=1$ its equations give
$X=-1/v$, $Y=u$, $Z=uv$, $u^2v^2=1$. It is the symmetry image
of the first and has the same rank and differential properties.

\subsection{Curvature when the signed coefficient is visible}

All statements in this subsection are over $\mathbb R$ with $d>0$.
We prove that every physical zero of the full parent, or of its
restriction to at most three hard matching events, has strictly
positive Hessian in its remaining angle and projective-state
coordinates. This supplies precisely the soft Hessian in (NC8).

\paragraph{No hard event.}
At a zero the transverse state block is positive by (NC6). If the
active exponents in (NC4) have rank one, stationarity (NC5) requires
an opposite pair on one of the four lines through the Newton square.
A horizontal opposite pair combines $\lambda D$ or $\alpha m$
with $\beta n$ or $cA$; in the four cases it also activates
$\lambda n,cD,\beta m,\alpha A$, respectively, off that line.
The vertical case is identical: combining $\alpha A$ or $\lambda n$
with $\beta m$ or $cD$ activates $\alpha m,cA,\beta n,\lambda D$.
The two diagonal pairs force exactly the supports
$\{n,m,A,D\}$ and $\{\lambda,c,\alpha,\beta\}$.
Any further event activates an exponent off the diagonal.

Neither diagonal support is physical. At $u=v=1$ the core matrix is
\[
 \begin{pmatrix}
 2&2&-2&-2\\2&6&-2&2\\-2&-2&6&-2\\-2&2&-2&2
 \end{pmatrix}.                                      \tag{NC12}
\]
The corresponding vectors $(1,0,0,1)$ and $(1,-1,1,1)$ are
annihilated by all selected events, have zero core value, and have
nonzero core image. A positive semidefinite matrix cannot have zero
quadratic value and nonzero image at the same vector.

If no outer exponent is active, the event support contains neither
neighboring nor opposite pairs in (NC10); these are exactly the
twelve activating pairs in (SG6). Such a set extends to a triple
with cyclic gaps $(2,3,3)$. Indeed an independent four-set in the
cycle would be alternating and contain an opposite pair, and the
only possible independent triple gap patterns are $(2,3,3)$ and
$(2,2,4)$, of which the latter contains an opposite pair. All the
allowed triples are cyclic images of $\{\lambda,m,c\}$.
The state
\[
 \psi_H=(uv,-u,u^2v,1+v^2+u^2v^2)
\]
annihilates this triple and has
$C(\psi_H)=-2u^2v^2D_uD_v<0$. This contradicts physicality.
Thus the active exponents span the plane.

\paragraph{One hard event: explicit coefficients.}
By symmetry take the hard event to be $\lambda$. Its kernel is
$W=-vX$. The constrained matrix in $(X,Y,Z)$ is positive definite
on $Z=0$, so it has two positive state pivots at a zero. Extracting
the coefficient of $\lambda$ from (NC2) gives its determinant as
$4D_uD_v^2UV L_\lambda(U,V)$, whose outer coefficients are
\[
\begin{array}{ll}
a_{00}=4\alpha t_7t_8,&
a_{01}=4(Dt_7t_6+2\alpha mt_8),\\
a_{02}=8mDt_6,&
a_{10}=4(2\alpha At_7+nt_8t_4),\\
a_{12}=8(\beta mt_6+cDt_2),&
a_{20}=8nAt_4,\\
a_{21}=8(\beta nt_2+cAt_4),&
a_{22}=8\beta ct_2.
\end{array}                                                    \tag{NC13}
\]
The factor follows directly by using first the normal coordinate
$W+vX$, whose packet is $4\lambda D_u(W+vX)^2$, and then the
three kernel coordinates. All $t_i$ in (NC13) are strictly positive.

The four rank-one possibilities are exhausted as follows. On the
horizontal line, off-line vanishing forces $\alpha=n=0$; its negative
end requires $D>0$, forcing $m=c=0$, so its positive end is absent.
On the vertical line, off-line vanishing gives $\alpha=D=0$; its
negative end requires $n>0$, forcing $\beta=0$, so its positive end
is absent. The slope-minus-one pair requires $m,D,n,A>0$, which
activates $a_{01}$ off the line. The slope-plus-one pair forces
$\alpha,\beta,c>0$ and $m=D=n=A=0$: it completes the perfect
matching containing the hard $\lambda$ event.

This last possibility is not nonnegative on the hard kernel. On the
matching curve every event and the core vanish, but (NC11) has a
nonzero derivative after only one event-normal direction is fixed.
All event squares have zero derivative there, so the constrained
parent changes sign. If no exponent is active, (NC13) gives
$\alpha=D=n=0$ and $\beta m=cA=\beta c=0$. Its support, including
the hard event, is contained in one of
\[
 \{\lambda,A,\beta\},\quad\{\lambda,m,c\},
 \quad\{\lambda,m,A\}.
\]
These are cyclic images of the negative-core triple above, so each
has a kernel witness contradicting constrained physicality. The
active exponents of $L_\lambda$ therefore span the plane.

\paragraph{Two hard events: exhaustive positivity table.}
Matching pairs have cyclic distances two, three or four, represented
by $\{\lambda,c\}$, $\{\lambda,m\}$ and
$\{\lambda,\alpha\}$. Extract the coefficient of the two hard
weights from (SG6). An adapted state basis shows that it is a
positive angle-dependent factor times the determinant of the
constrained two-by-two state form. After the monomial $UV$ is removed,
the exact positivity conditions of its eight outer coefficients are
the following. Juxtaposition means both entries are positive;
``or'' means at least one condition holds.
\[
\begin{array}{c|c|c|c}
 e&\{\lambda,m\}&\{\lambda,c\}&\{\lambda,\alpha\}\\\hline
(-1,-1)&\alpha&\alpha&\text{always}\\
(-1,0)&D\text{ or }\alpha&D\text{ or }\alpha m&D\text{ or }m\\
(-1,1)&D&Dm&Dm\\
(0,-1)&n\text{ or }A\alpha&n&A\text{ or }n\\
(0,1)&\beta\text{ or }Dc&D\text{ or }\beta m&Dc\text{ or }\beta m\\
(1,-1)&An&0&An\\
(1,0)&Ac\text{ or }\beta n&A&Ac\text{ or }\beta n\\
(1,1)&\beta c&\beta&\beta c
\end{array}                                                    \tag{NC14}
\]
This table is obtained by differentiating (SG6), not by a sign
approximation. For example the first two entries in the
$\{\lambda,m\}$ column, apart from the positive factor four, are
$2\alpha(2A+2D+2c+d)$ and
$2(2AD+2A\alpha+4D\alpha+2Dc+Dd+2\alpha c+\alpha d)$.

For $\{\lambda,m\}$, either diagonal pair activates the off-line
exponent $(-1,0)$. A horizontal pair combines $D$ or $\alpha$ with
$Ac$ or $\beta n$ and activates an off-line exponent in all four
cases. A vertical pair combines $n$ or $A\alpha$ with $\beta$ or
$Dc$ and does the same. Rank zero forces
$\alpha=D=n=\beta=0$ and $Ac=0$, so only $A$ or only $c$ may remain.

For $\{\lambda,c\}$ the slope-minus-one pair is impossible because
one end is identically zero. A horizontal pair requires $A>0$ and
$D>0$ or $\alpha m>0$, activating $(0,1)$ or $(-1,-1)$ off the
line. A vertical pair requires $n>0$ and either $D>0$ or
$\beta m>0$, activating $(-1,0)$ or $(1,1)$ off the line.
The slope-plus-one pair forces $\alpha,\beta>0$ and
$D=n=A=m=0$. Rank zero leaves only $m$.

For $\{\lambda,\alpha\}$ the exponent $(-1,-1)$ is always active.
A rank-one stationary support must be the slope-plus-one pair,
forcing $c,\beta>0$ and $D=m=A=n=0$. These last two rank-one
possibilities complete $\{\lambda,c,\alpha,\beta\}$. Imposing
only two hard events still leaves a nonzero core derivative in
(NC11), so the constrained parent changes sign on the matching
curve. The rank-zero possibilities are contained in
$\{\lambda,m,A\}$ or $\{\lambda,m,c\}$. At $u=v=1$, their kernel
vectors $(1,-1,0,2)$ and $(1,-1,1,3)$ have core values $-4,-8$.
Thus every obstruction is excluded by constrained physicality.
The active exponents again span the plane.

\paragraph{Three hard events: two complete Laurent restrictions.}
A three-edge matching has cyclic gaps $(2,2,4)$ or $(2,3,3)$;
rotations and reflections are transitive on each pattern. The two
orbits have representatives and kernel states
\[
 M_N=\{\lambda,m,A\},\quad\psi_N=(uv,-u,0,1+v^2);
 \qquad
 M_E=\{n,m,A\},\quad\psi_E=(u,0,0,v).
\]
Direct substitution into the entire parent gives
\begin{align}
 {F(\psi_N)\over2UV(1+V)}
 ={}&{2n(1+U)\over V}+2cU+{2D(1+V)\over U}
       +{2\alpha(1+U)(1+V)\over UV}+2\beta V-d,
                                                        \tag{NC15}\\
 {F(\psi_E)\over2UV}
 ={}&d(U/V-1)+{2\lambda(1+U)\over V}
       +{2cU(1+V)\over V}+{2D(U-V)^2\over UV}
       +{2\alpha(1+U)\over U}+2\beta(1+V).            \tag{NC16}
\end{align}
For (NC15), any positive $n,D$ or $\alpha$ supplies two independent
exponents. If they vanish, $2cU+2\beta V-d$ cannot have a stationary
zero: stationarity forces $c=\beta=0$ and the value is $-d<0$.
For (NC16), $d>0$ supplies $(1,-1)$. Rank less than two forces
$\lambda=c=\alpha=\beta=0$. The remaining function of $x=U/V$
has infimum
\[
 2\sqrt{2D(d+2D)}-d-4D<0,
 \qquad (d+4D)^2-8D(d+2D)=d^2>0.
\]
The formula also gives a negative infimum when $D=0$. This contradicts
physicality. Both constrained Laurent Hessians are therefore strict
at every zero.

\paragraph{From angular determinants to the full soft Hessian.}
For zero, one or two hard events, (NC6) supplies a positive state
block on the remaining projective-state coordinates. In a fixed real
kernel chart write the constrained state form as
$a+2b^Ty+y^TC_0y$, with $C_0\succ0$ at contact. Its state Schur
complement is $a-b^TC_0^{-1}b$, and its determinant equals this
scalar times $\det C_0$. At a stationary zero, differentiating a
positive prefactor does not change the Hessian sign. Equations
(NC5), (NC12)--(NC14) therefore give a positive angular Schur
complement and hence the full positive soft Hessian. For three hard
events only the two angles remain and (NC15)--(NC16) apply directly.
All these matrix inverses are ordinary real residue calculations.
The exact certificate itself is (NC7)--(NC9), with no tensor inverse.

\subsection{Four matching weights, with no comparability assumption}

Suppose the four positive weights of a perfect matching dominate
$d$ and every complementary event. By symmetry use
$\{n,m,A,D\}$ and arrange that $D=\mu=\min(n,m,A,D)$.
Set
\[
 \eta_1=\mu/n,\quad\eta_2=\mu/A,\quad\eta_3=\mu/m,
 \qquad0<\eta_i\le1.
\]
Their residues may vanish independently. Work near the matching
curve in $W=1$ and introduce the exact state coordinates
\[
 R=v-X+uvY-uZ,\qquad
 Z={v-X+uvY-R\over u},\qquad\delta=v-u,
\]
using only the testing unit $u^{-1}$. The fourth event is exactly
$f+L$, where
\[
 f={u^2-v^2\over u}=-2\delta-\delta^2/u,\qquad
 L=v(u+u^{-1})X-(1+v^2)Y+(v/u)R.                    \tag{NC17}
\]
The first three hard packets are
$4nD_uD_vX^2+4AD_uD_vY^2+4mR^2$.
On $X=Y=R=0$ the core is exactly $2uf$, so
\[
 4\mu f^2+2duf
 =4\mu\left(f+{du\over4\mu}\right)^2
                         -{d^2u^2\over4\mu}.         \tag{NC18}
\]
Let $\tau=\max(\lambda,c,\alpha,\beta)$. Testing (NC18) at
$u=1$, $v=\sqrt{1+d/(4\mu)}$ proves
\[
 \tau>0,\qquad \rho={d^2\over\mu\tau}\text{ bounded}.
 \tag{NC19}
\]
Indeed all complementary amplitudes are bounded at this source-field
test; if their scale were infinitesimal relative to $d^2/\mu$, the
complete parent would be negative. The case $d=0$ was already an
event-square sum. Put $\varepsilon=d/\mu$, which is infinitesimal.

If $\bar\rho=0$, use the five rows
\[
 (\sqrt{n/\tau}X,\sqrt{A/\tau}Y,
       \sqrt{m/\tau}R,\sqrt{\mu/\tau}\delta,1).
\]
Their four-normal residue is, with $B_0=1+u_0^2$,
\[
4B_0^2z_1^2+4B_0^2z_2^2+4z_3^2
 +4(-2z_4+B_0\sqrt{\bar\eta_1}z_1
       -B_0\sqrt{\bar\eta_2}z_2+\sqrt{\bar\eta_3}z_3)^2,
 \tag{NC20}
\]
which is positive definite. The constant residue is
$4(1+u_0^2)(P+Qu_0^2)>0$, where
$P=\overline{(\lambda+\alpha)/\tau}$,
$Q=\overline{(c+\beta)/\tau}$ and $P+Q>0$.
The signed linear coefficients are $\sqrt{\rho\eta_i}$ and
$\sqrt\rho$, so vanish. All other normalized coefficients are
bounded by $\varepsilon,\eta_i,\tau/\mu$ and bounded testing
units. This is an exact strict bounded Gram of the whole parent.

Now suppose $\bar\rho>0$. In the source-field stationary equations
put
\[
 X=\varepsilon\eta_1s_1,\quad
 Y=\varepsilon\eta_2s_2,\quad
 R=\varepsilon\eta_3s_3,\quad\delta=\varepsilon r,
 \quad u=t.                                          \tag{NC21}
\]
Divide the three state equations and the $\delta$ equation by $d$,
and the tangent equation by $\tau$. All coefficients are bounded
after clearing powers of the unit $t$: the factors $\eta_i$ cancel
the large weights $\mu/\eta_i$, and
$\tau=\mu\varepsilon^2/\rho$. Put, in residue,
\[
 B_t=1+t^2,\quad\ell=(B_t,-B_t,1),\quad
 b=(B_t^2,B_t^2,1),\quad
 k=-2r+\sum_i\bar\eta_i\ell_i s_i.
\]
Direct differentiation of (NC1), (NC17) gives the four normal equations
\[
 8b_i s_i+8k\ell_i+(4tB_t,-4tB_t,0)_i=0,
 \qquad-16k-4t=0.                                   \tag{NC22}
\]
Eliminating the last row from the first three gives diagonal pivots
$8b_i$, with the last pivot $32$. The determinant is
$32\cdot8^3B_t^4>0$, even if some $\bar\eta_i=0$.
Their solution is
\[
 s_1^*=-{t\over4B_t},\quad s_2^*={t\over4B_t},\quad
 s_3^*={t\over4},\quad
 r^*={t\over8}(1-\bar\eta_1-\bar\eta_2+\bar\eta_3).
 \tag{NC23}
\]
The complete normalized residue objective is
\[
\begin{split}
 \Psi={}&\bar\rho\{4\sum_i\bar\eta_i b_i s_i^2+4k^2
           +4tB_t(\bar\eta_1s_1-\bar\eta_2s_2)-4tr\}\\
 &+4(1+t^2)(P+Qt^2).
\end{split}
\]
Completing its four normal squares by (NC23) leaves exactly
\[
 g(t)=4(1+t^2)(P+Qt^2)
       -{\bar\rho t^2\over4}
          (1+\bar\eta_1+\bar\eta_2+\bar\eta_3).     \tag{NC24}
\]
This completion is polynomial in the $\bar\eta_i$ and remains
valid when one vanishes; its scaled equation in (NC22) still has
the displayed positive pivot. Substituting the corresponding normal
values into the physical source and taking residue proves $g(t)\ge0$
for every real $t\ne0$.

If $g(t_0)>0$, shifting by (NC23) gives the positive normal form
(NC20) and a positive constant, hence a strict bounded Gram. At a
zero, divide (NC24) by $4t^2$:
\[
 P(1+t^{-2})+Q(1+t^2)
       -{\bar\rho\over16}(1+\bar\eta_1+\bar\eta_2+\bar\eta_3).
\]
Stationarity gives $P/t_0^2=Qt_0^2$. Since $P+Q>0$, both are
positive; the logarithmic second derivative is positive, and thus
$g''(t_0)>0$. Append the tangent equation to (NC22). Its Schur
complement is $g''(t_0)$, so its real Jacobian is invertible. The
semialgebraic implicit-function argument gives an exact source
stationary center of the form (NC21), with bounded $s_c,r_c,t_c$.
Its normalized value is nonnegative by physicality.

The exact full Taylor rows are
\[
\big(\sqrt{n/\tau}(X-X_c),\sqrt{A/\tau}(Y-Y_c),
\sqrt{m/\tau}(R-R_c),\sqrt{\mu/\tau}(\delta-\delta_c),t-t_c\big).
 \tag{NC25}
\]
Here is the required coefficient control, including all normal terms.
For $a_i=\mu/\eta_i$ the two center factors in a hard packet are
\[
 {a_ix_{ci}^2\over\tau}=\rho\eta_i s_{ci}^2,
 \qquad {a_ix_{ci}\over\sqrt{a_i\tau}}
                    =\sqrt{\rho\eta_i}\,s_{ci}.      \tag{NC26}
\]
The fourth packet has the analogous factors $\rho r_c^2$ and
$\sqrt\rho r_c$. The core belongs to the ideal $(X,Y,R,\delta)$
because it vanishes identically on the matching curve. Its center-linear
factors after division by $\tau$ are therefore
$d x_{ci}/\tau=\rho\eta_i s_{ci}$ and
$d\delta_c/\tau=\rho r_c$; its normalized normal factors are
$d/\sqrt{a_i\tau}=\sqrt{\rho\eta_i}$ and
$d/\sqrt{\mu\tau}=\sqrt\rho$. Its quadratic normal factors are
bounded by $d/\mu=\varepsilon$. Complementary coefficients are at
most $\tau$ and give smaller factors. Each packet is at most quadratic in the
state coordinates, and all remaining higher Taylor factors are
actual angle or state deviations with zero residue. The only testing
denominators are powers of the unit $t$. Thus exact packetwise Taylor
allocation is a bounded Gram on (NC25). Its normal residue is (NC20)
and its tangent Schur complement is $g''(t_0)/2>0$, as follows by
differentiating the completed squares (NC23)--(NC24). This assertion
also holds when $\bar\eta_i=0$: the normal form (NC20), its cross
entries and the completed Schur expression are continuous polynomial
expressions in $\sqrt{\bar\eta_i}$, with a strict normal block.
The strict-Gram lemma, the nonnegative center value and pure-factor
restoration prove the four-event case without any restriction
on its internal coefficient ratios.

\subsection{A terminating induction on event constraints}

Normalize the nonzero source by $d+\lambda+n+m+c+A+D+\alpha+\beta=1$.
The factor is positive and pure-source. Literal $d=0$ is already an
event-square sum; positive slack residue is covered by the positive
residue lemma. If $\bar d>0$ and the residue slack is zero, the
unconstrained curvature argument and (NC9) with $k=0$ apply.

It remains that $d>0$, $\bar d=0$ and the residue slack is zero.
Product residue states were settled above. At an entangled residue,
let $S$ be all events with positive coefficient residue. This is a
nonempty matching. If it has four edges, every complementary
coefficient and $d$ is infinitesimal relative to its smallest weight,
and (NC17)--(NC26) apply. Otherwise use its at most three independent
event amplitudes as the exact hard state coordinates. Set
\[
 \sigma=\max\bigl(\{d\}\cup\{c_e:e\notin S\}\bigr)>0.
\]
Every hard weight $a_i$ satisfies $\sigma/a_i$ infinitesimal. The
parent has the diagonal form (NC7), and its constrained next-scale
residue $g=\overline G(0,y)$ is nonnegative by physicality.

If $g(y_0)>0$, (NC9)'s positive-value argument finishes. If $g(y_0)=0$
and $\overline{d/\sigma}>0$, the one-, two- or three-hard curvature
calculation above gives a strict full soft Hessian; (NC8)--(NC9)
finish without assuming comparability of the hard weights.

The only remaining case has $g(y_0)=0$ and
$\overline{d/\sigma}=0$. Then $g$ is a nonnegative event-square
sum, and at least one remaining event coefficient has positive
normalized residue. All such event amplitudes vanish at the same
entangled residue state. Add them all to $S$. The enlarged set is
still a matching; every coefficient left outside it is infinitesimal
relative to each of its weights. If it has at most three edges, repeat.
If it has four, apply the four-event proof above, which allows arbitrary source scales.
Each repetition adds an edge, and a matching has at most four edges.
There are therefore at most three enlargements. Every step retains
the entire original parent. This proves
Theorem~\ref{thm:noncommuting-refined}.



